\section{序論}
% 近年，少子高齢化の進行に伴い，ロボットが人間の労働を代替する役割への期待が高まっている．
% しかし，現状の産業用ロボットは動作を実行するために人間によるプログラミングを必要とし，環境の変化に柔軟に適応することが難しい．
% このような課題に対応するため，近年では機械学習をロボットに応用する研究が進められている．機械学習モデルを用いることで，ロボットがデータから動作方策を学習し，環境変化にロバストな動作生成が可能となる手法が注目されている．

近年，機械学習をロボットに応用する研究が進められている．
その中でも遠隔操作技術の一種であるバイラテラル制御を模倣学習に用いる手法は力加減を含めた動作の模倣が可能であり，多くのタスクに対してその有効性が示されている\cite{Sasagawa_2020,Pancake-Saigusa,Akagawa202322002155,Kobayashi202424004380}
この手法ではタスクに成功したデータから学習するため，ロボットが試行錯誤的に環境と相互作用して学習する強化学習\cite{ReinforcementLearning-Han}と比較してサンプル効率が高く，また人間が教示した動作データから End-to-End で学習を行うため，プログラミングによる手法よりも運用が容易である．

% ここの文章意味不明では？じゃあ外部センサで見ろよという話になる
しかし，模倣学習には教師データと大きく異なる状況においては適応が困難な場合がある．
例えば学習済みのロボットの動作経路上に障害物が出現したとき，ロボットはその障害物を回避できない．
このような場合，環境が変わるたびにデータの取得と再学習を行う必要がある．

模倣学習においてロボットが環境を解釈し，適切な動作を計画するための手法として，学習した軌道を感覚フィードバックによって修正する方法がある．
Ashara らは，動的動作プリミティブ(DMP)を用いた模倣学習において，人間の操作データから様々な種類の障害物回避行動を学習しモデル化する手法を提案した\cite{Learning-baseAvoidanceWithAPF-Rai}．
また Ginesiらは，障害物の形状や相対運動に基づくポテンシャル関数を DMP に適用し，様々な静的環境において障害物回避動作動作を達成している\cite{FixedParameterAvoidanceWithAPF-Ginesi}．
しかしこれらの手法は力情報を用いておらず，環境との接触を含むタスクへの適用が困難であると考えられる．
また，これらのようなリアクティブな障害物回避手法は，未だバイラテラル制御による模倣学習に適用された例はない．

そこで，バイラテラル制御を用いた模倣学習を行ったモデルに対して人工ポテンシャル場法を適用することを提案した．
人工ポテンシャル場法では，ポテンシャル場という概念を用いて，目標への吸引力や障害物からの斥力をロボットに適用し，その合力に基づいて移動方向を決定する\cite{APF-Khatib1985}．
これによって自律動作中のロボットの入力に仮想的な力を加えて軌道を変化させ，必要な学習データを増やすことなく障害物を回避する動作を実現する．

しかし人工ポテンシャル場法を適用して自律動作を行う際には，学習時に存在しない力をロボットに加えることになる．
従来模倣学習で用いられてきた Long Short-Term Memory (LSTM)\cite{LSTM} では，自律動作時の入力が学習データと乖離した場合の予測が難しく，ポテンシャル場法を適用することが困難であると考えられる．
これは，LSTM が時系列の関係を学習する際に過去の情報を保持するため，内部状態を持つという特性による．

そこで，Hayashiらの提案した階層型モデル\cite{HierarchicalModel-Hayashi}を拡張したモデルを提案した．
階層型モデルでは上位層がロボットの数ステップ先の状態を予測し，下位層が現在の状態と上位層が予測した将来の状態の間の動作を予測する．
Hayashiらによって提案されたモデルでは上位層と下位層がそれぞれ LSTM で構成されている．
これに対し本研究では，上位層は内部状態を持つ LSTM，下位層は内部状態を持たない Multilayer perceptron (MLP)\cite{MLP}とした．
これにより上位層のみが時系列の関係を学習し抽象度の高い予測を行い，下位層は現在の状態と上位層の予測を補完するような予測を行うため，ポテンシャル場に対して適応的な動作が可能になると考えられる．

% などの内部状態を持つNN
% また模倣学習においては動作の時系列の関係が重要であり，このため学習に使用するモデルには Long Short-Term Memory (LSTM)\cite{LSTM}などの内部状態を持つNNが用いられることが多い．
% しかし本研究では自律動作時に学習時に存在しない力をロボットに加えることになるため， LSTM などの内部状態を持つ NN を用いたとき，その非マルコフ性により予期せぬ動作を行う場合がある．
% そこで本研究では学習モデルとして，Hayashiらの提案した階層型モデル\cite{HierarchicalModel-Hayashi}を拡張したモデルを用いた．
% 階層型モデルでは上位層がフォロワの数ステップ先の状態を予測し，下位層が現在の状態と上位層が予測した将来の状態の間の動作を予測する．ここで上位層はシステムの非マルコフ的な性質を扱っており，下位層は，既知の情報である現在の状態と予測された将来の状態の間の関係を補完すればよいため，マルコフ性を持つといえる．
% したがって本研究で用いる学習モデルでは，上位層は内部状態を持つ LSTM ，下位層は内部状態を持たない Multilayer perceptron (MLP)\cite{MLP}とした．
% これにより上位層のみが時系列の関係を学習するため，より抽象度の高い予測を行い，ポテンシャル場に対して適応的な動作が可能になると考えられる．

本研究では，バイラテラル制御を用いた模倣学習の自律動作時にポテンシャル場を適用することで環境の変化に対して適応的な動作が可能であるかを Pick-and-Place タスクによって検証した．また提案手法における階層型学習モデルの有効性を， LSTM のみで構成された非階層型モデルと比較することで検証した．